% 第 1 章：架构总览与方法族
\section{架构总览与方法族}\label{sec:rel-overview}

\subsection{功能定位}
可靠性评估模块回答“系统在元件随机故障下少供多少电、失负荷多频繁多久”。全部方法
遵循同一建模主线（\srcpath{docs/theory/reliability\_mathematical\_models\_and\_intelligent\_cyber\_physical\_assessment.md}
第 1 节的中心原理）：
\begin{equation}
\text{可靠性}=\text{事件频率}\times\text{类条件后果}\times\text{持续时间}.
\end{equation}
命名空间为 \texttt{hacdcpf::analysis}。

\subsection{统一数据流}
\begin{figure}[H]\centering
\begin{tikzpicture}[font=\footnotesize,
  box/.style={draw,rounded corners,minimum height=1.0cm,text width=2.7cm,align=center,inner sep=4pt},
  arr/.style={->,thick}]
  \node[box] (rich) at (0,0) {富语义模型\\\texttt{HybridPowerSystem}};
  \node[box] (res) at (3.7,0) {参数解析\\$\lambda,r,U$ 规范化};
  \node[box] (enum) at (7.4,0) {状态/故障枚举\\抽样·N-1·目录};
  \node[box] (patch) at (11.1,0) {后果算子\\$G_m{=}\Phi_m(G_0,x_m)$};
  \node[box] (shed) at (3.7,-2.5) {切负荷引擎\\DC-OPF·LP·MILP};
  \node[box] (metric) at (7.4,-2.5) {指标聚合\\EENS·LOLE·LOLF};
  \node[box] (audit) at (11.1,-2.5) {口径审计\\\fld{model\_scope}};
  \draw[arr] (rich)--(res); \draw[arr] (res)--(enum); \draw[arr] (enum)--(patch);
  \draw[arr] (patch.south) |- (shed.west);
  \draw[arr] (shed)--(metric); \draw[arr] (metric)--(audit);
\end{tikzpicture}
\caption{可靠性评估统一数据流：解析$\to$枚举$\to$后果映射$\to$切负荷求解$\to$指标$\to$口径审计。}
\end{figure}

\subsection{公共入口族}
全部入口位于命名空间 \texttt{hacdcpf::analysis}；头文件除三阶段模型在
\srcpath{include/hacdcpf/analysis/} 外，均在 \srcpath{include/hacdcpf/reliability/}，下表“头文件”列仅
记基名，函数名与“入口”列同名：
\begin{center}\footnotesize
\begin{tabular}{@{}L{4.9cm}L{4.9cm}L{3.8cm}@{}}
\toprule
入口（函数） & 头文件 & 方法族\\
\midrule
\srcpath{resolve\_reliability\_params} & \srcpath{reliability\_assessment.hpp} & 参数解析（全方法共用）\\
\srcpath{run\_nonsequential\_mc} & \srcpath{reliability\_assessment.hpp} & 非序贯 MC\\
\srcpath{run\_sequential\_mc} & \srcpath{reliability\_assessment.hpp} & 序贯 MC（含 LOLF）\\
\srcpath{run\_distribution\_fmea} & \srcpath{reliability\_assessment.hpp} & 元件级 FMEA\\
\srcpath{run\_failure\_mode\_fmea} & \srcpath{failure\_mode.hpp} & 故障模式级 FMEA（N-2）\\
\srcpath{run\_three\_stage\_reliability} & \srcpath{three\_stage\_reliability.hpp} & 三阶段恢复 MILP\\
\srcpath{run\_frequency\_duration\_analysis} & \srcpath{reliability\_assessment.hpp} & 频率-持续时间解析\\
\srcpath{build\_failure\_mode\_catalog} & \srcpath{failure\_mode.hpp} & 故障模式目录\\
\srcpath{evaluate\_failed\_network\_state} & \srcpath{reliability\_assessment.hpp} & 后果引擎入口\\
\srcpath{compute\_distribution\_indices} & \srcpath{reliability\_assessment.hpp} & SAIFI/SAIDI/ASAI\\
\bottomrule
\end{tabular}
\end{center}

\subsection{物理后果引擎阶梯}
不同方法调用不同保真度的切负荷引擎（\srcpath{docs/theory/reliability\_mathematical\_models\_and\_intelligent\_cyber\_physical\_assessment.md}
第 4.3 节）：
\begin{center}\footnotesize
\begin{tabular}{@{}L{3.6cm}L{4.2cm}L{6.4cm}@{}}
\toprule
引擎 & 使用者 & 有效边界\\
\midrule
AC DC-OPF & 纯交流 MC/FMEA & 无电压幅值/无功可行性（$P=B\theta$）\\
混合 AC/DC 网络 LP & 混合 MC、元件/故障模式 FMEA & DC 传输线性稳态；无非线性 AC 校验\\
耦合 LinDistFlow 恢复 MILP & 三阶段模型 & 线性化损耗；拒绝 LCC 与多端口能量路由器\\
\bottomrule
\end{tabular}
\end{center}

\subsection{诚实结果口径}
每个结果对象都暴露 \fld{model\_scope}、\fld{model\_limitations} 与 \fld{ValidityFlags}，
是数学结果的组成部分，下游 API 必须保留：
\begin{itemize}
  \item MC/FMEA 在纯交流上为 \texttt{"ac-only-dcopf"}，
        \fld{validity.\allowbreak dc\_load\_curtailment\_included}=false——EENS/LOLE 不含 DC 负荷中断；
  \item 混合系统 FMEA 升级为 \texttt{"hybrid-acdc-network-lp"}；
  \item 三阶段模型为
        \texttt{"coupled-acdc-\allowbreak lindistflow-\allowbreak restoration-milp"}，
        含 LCC/能量路由器
        时清 \fld{ok} 并置物理有效性标志为假；
  \item \fld{critical\_components.importance} 是共现归因（可相加超 100\%），\textbf{非} Birnbaum
        边际重要度（\srcpath{reliability\_assessment.hpp} \fld{ReliabilityResult} 注释，:ReliabilityResult）。
\end{itemize}
